% Pista geo-24i-q6 · Geometria
% Procedència: PAU Catalunya 2024 · Convocatòria de juny (incidències) · Sèrie 5 · Qüestió 6
\documentclass[11pt,a4paper]{article}
\usepackage{pista}
\pistainfo{Catalunya 2024}{Convocatòria de juny (incidències)}{Sèrie 5}

\begin{document}

\section*{\textcolor{blauPista}{Qüestió 6}}

\noindent Considereu les rectes
\[
r: \dfrac{x-5}{4} = \dfrac{y-4}{3} = \dfrac{z-3}{-1} \quad\text{i}\quad s: \begin{cases} x = 4 + 2k \\ y = 3 + k \\ z = -1 \end{cases}
\]

\subsection*{\textit{a)} \normalfont Quina és la seva posició relativa? Calculeu l'equació implícita d'un pla $\pi$ que sigui paral·lel a les dues rectes i que passi per l'origen de coordenades. \hspace*{\fill}\textnormal{[1,25~punts]}}

\pista{Dedueix els vectors directors $\vec{v}_{r}$ i $\vec{v}_{s}$ i un punt de cada recta ($P \in r$, $Q \in s$).

\medskip
\textit{Posició relativa.}
\begin{itemize}
\item Si els directors són proporcionals, les rectes són paral·leles (o coincidents).
\item Si no ho són, es tallen o s'encreuen. Per a distingir-ho, calcula el determinant format per $\vec{v}_{r}$, $\vec{v}_{s}$ i $\vec{PQ}$.
\end{itemize}
\textit{El pla $\pi$.}
\begin{itemize}
\item El vector normal ha de ser perpendicular als dos directors:
\[
\vec{n}_{\pi} = \vec{v}_{r} \times \vec{v}_{s}
\]
\item Com que el pla passa per l'origen, el terme independent de l'equació general és $D = 0$.
\end{itemize}}

\subsection*{\textit{b)} \normalfont Calculeu l'equació de la recta $t$ que talla les dues rectes $r$ i $s$ perpendicularment. \hspace*{\fill}\textnormal{[1,25~punts]}}

\pista{La recta $t$ és la «perpendicular comuna».
\begin{itemize}
\item \textit{Direcció:} $\vec{v}_{r} \times \vec{v}_{s}$, que ja has calculat a l'apartat anterior.
\item \textit{Un punt:} pren un punt genèric de $r$ (paràmetre $\lambda$) i un de $s$ (paràmetre $k$), i imposa que el vector que els uneix sigui perpendicular alhora a $\vec{v}_{r}$ i a $\vec{v}_{s}$. Obtindràs un sistema de dues equacions amb incògnites $\lambda$ i $k$.
\item Substitueix el valor d'un dels paràmetres a l'equació de la seva recta: tindràs un punt concret de $t$.
\end{itemize}
Amb el punt i la direcció, ja pots escriure l'equació que busques.}

\end{document}
